% =====================================================================================================
% ĐỀ XUẤT: FISA trung bình nhân có trọng số bộ ba (FISA-GM-w) — đưa vào Chương 3, ngay sau định nghĩa FISA-GM.
% Khớp mã: kfold/weighted.py (triple_weights), kfold/core.py (fisa_gm, weights=...), cấu hình config_ch3.json.
% Ký hiệu cần đối chiếu với chương: r (số thuộc tính), \mathcal{T} (tập bộ ba), C_l(t;x) (FISA), \varepsilon.
% =====================================================================================================

\begin{dn}[Điểm bằng chứng của bộ ba]\label{dn:mm.triplescore}
Cho FKG xây từ cơ sở luật huấn luyện, $\mathcal{T}=\{t=(a,b,c): 1\le a<b<c\le r\}$ là tập các bộ ba thuộc tính,
$|\mathcal{T}|=\binom{r}{3}$, và $C_l(t;\mathbf{x})$ là tổng trọng số của các luật nhãn $l$ khớp truy vấn
$\mathbf{x}$ trên bộ ba $t$ (Định nghĩa~\ref{dn:fisa}). Với bài toán hai lớp, \emph{điểm bằng chứng} của bộ ba
$t$ cho truy vấn $\mathbf{x}$ là
\begin{equation}\label{eq:mm.triplescore}
  s_t(\mathbf{x}) = \log\bigl(C_1(t;\mathbf{x})+\varepsilon\bigr) - \log\bigl(C_0(t;\mathbf{x})+\varepsilon\bigr),
  \qquad \varepsilon>0 .
\end{equation}
FISA-GM tương ứng với điểm $s(\mathbf{x})=\frac{1}{|\mathcal{T}|}\sum_{t\in\mathcal{T}} s_t(\mathbf{x})$.
\end{dn}

\begin{dn}[Trọng số phân biệt ước lượng chéo]\label{dn:mm.tripleweight}
Gọi $\mathcal{D}^{0}=\{(\mathbf{x}_i,y_i,g_i)\}_{i=1}^{n_0}$ là các \emph{ảnh gốc} của tập huấn luyện ($g_i$ là mã
bệnh nhân; không gồm mẫu sinh bởi SMOTE, không gồm tập xác thực và kiểm tra), đã mờ hoá bằng chính phép mờ hoá
của mô hình. Chia $\mathcal{D}^{0}$ thành $K$ phần $P_1,\dots,P_K$ phân tầng theo nhãn sao cho mỗi bệnh nhân
thuộc đúng một phần. Với mỗi $k$, xây FKG$^{(-k)}$ từ $\mathcal{D}^{0}\setminus P_k$ và tính điểm ngoài phần
$\tilde s_t(\mathbf{x}_i)=s^{(-k)}_t(\mathbf{x}_i)$ với $i\in P_k$. Trọng số của bộ ba $t$ là
\begin{equation}\label{eq:mm.tripleweight}
  w_t = \max\Bigl(\operatorname{AUC}\bigl(\{\tilde s_t(\mathbf{x}_i), y_i\}_{i=1}^{n_0}\bigr) - \tfrac12,\; 0\Bigr),
  \qquad
  \hat w_t = \frac{w_t}{\sum_{u\in\mathcal{T}} w_u},
\end{equation}
trong đó $\operatorname{AUC}$ là diện tích dưới đường ROC (xác suất một ảnh dương có điểm cao hơn một ảnh âm,
tính hoà bằng $\tfrac12$). Nếu $\sum_u w_u=0$, đặt $\hat w_t = 1/|\mathcal{T}|$.
\end{dn}

\begin{dn}[FISA trung bình nhân có trọng số — FISA-GM-w]\label{dn:mm.fisagmw}
Với trọng số $\hat{\mathbf w}$ của Định nghĩa~\ref{dn:mm.tripleweight} và FKG xây từ \emph{toàn bộ} cơ sở luật
huấn luyện, giá trị quyết định của lớp $l$ cho truy vấn $\mathbf{x}$ là trung bình nhân có trọng số
\begin{equation}\label{eq:mm.fisagmw}
  D^{\mathrm w}_l(\mathbf{x}) = \prod_{t\in\mathcal{T}} \bigl(C_l(t;\mathbf{x})+\varepsilon\bigr)^{\hat w_t}
  = \exp\Bigl(\sum_{t\in\mathcal{T}} \hat w_t \log\bigl(C_l(t;\mathbf{x})+\varepsilon\bigr)\Bigr),
\end{equation}
nhãn dự đoán là $\hat y(\mathbf{x})=\arg\max_l D^{\mathrm w}_l(\mathbf{x})$, và với hai lớp, điểm xếp hạng là
\begin{equation}\label{eq:mm.fisagmw.score}
  s_{\mathrm w}(\mathbf{x}) = \log D^{\mathrm w}_1(\mathbf{x}) - \log D^{\mathrm w}_0(\mathbf{x})
  = \sum_{t\in\mathcal{T}} \hat w_t\, s_t(\mathbf{x}),
  \qquad \hat y(\mathbf{x}) = \mathbb{1}\bigl[s_{\mathrm w}(\mathbf{x})>0\bigr].
\end{equation}
\end{dn}

\begin{md}[Tính chất của FISA-GM-w]\label{md:mm.fisagmw}
\begin{enumerate}[(i)]
  \item \emph{Tổng quát hoá FISA-GM:} nếu $w_t$ bằng nhau với mọi $t$ thì $D^{\mathrm w}_l = D^{\mathrm{GM}}_l$,
        do đó nhãn và điểm trùng FISA-GM.
  \item \emph{Bị chặn:} với mọi $l$ và $\mathbf{x}$,
        $\min_{t:\hat w_t>0} \bigl(C_l(t;\mathbf{x})+\varepsilon\bigr) \le D^{\mathrm w}_l(\mathbf{x})
        \le \max_{t:\hat w_t>0} \bigl(C_l(t;\mathbf{x})+\varepsilon\bigr)$.
  \item \emph{Loại bộ ba không phân biệt:} bộ ba có $\operatorname{AUC}_t\le\tfrac12$ có $\hat w_t=0$ và không ảnh
        hưởng tới quyết định. Bộ ba chỉ gồm các thuộc tính độc lập với nhãn có $\mathbb{E}[\operatorname{AUC}_t]=\tfrac12$
        nên nhận trọng số gần $0$ khi $n_0$ đủ lớn; bộ ba trộn thuộc tính nhiễu với thuộc tính có thông tin được
        trọng số theo mức phân biệt còn lại. Ở FISA-GM, mọi bộ ba có cùng trọng số $1/|\mathcal{T}|$.
  \item \emph{Không rò rỉ:} $\hat{\mathbf w}$ chỉ phụ thuộc vào ảnh gốc của tập huấn luyện; nhãn của tập xác
        thực, tập kiểm tra và các mẫu SMOTE không tham gia ước lượng.
  \item \emph{Bất biến thứ tự:} $\operatorname{AUC}_t$ chỉ phụ thuộc thứ tự của $\tilde s_t$, nên $w_t$ không đổi
        khi $\tilde s_t$ bị biến đổi tăng ngặt (ví dụ tịnh tiến bởi một hằng số).
\end{enumerate}
\end{md}
\noindent\emph{Chứng minh.} (i) $\hat w_t=1/|\mathcal{T}|$ cho đúng \eqref{eq:mm.fisagmw} của FISA-GM.
(ii) $\log D^{\mathrm w}_l$ là tổ hợp lồi của $\log(C_l(t;\mathbf{x})+\varepsilon)$ trên các $t$ có $\hat w_t>0$.
(iii) Phần đầu trực tiếp từ \eqref{eq:mm.tripleweight} và \eqref{eq:mm.fisagmw.score}; nếu mọi thuộc tính của
$t$ độc lập với nhãn thì $\tilde s_t(\mathbf{x}_i)$ độc lập với $y_i$, nên $\mathbb{E}[\operatorname{AUC}_t]=\tfrac12$.
(iv) Theo xây dựng của Định nghĩa~\ref{dn:mm.tripleweight}.
(v) AUC bất biến với mọi biến đổi tăng ngặt của điểm. \hfill$\square$

\begin{nx}[Lựa chọn thiết kế]\label{nx:mm.fisagmw}
\begin{itemize}
  \item Trọng số dùng $\operatorname{AUC}_t-\tfrac12$ (không âm) thay vì hệ số hồi quy: không cần chọn siêu
        tham số, không đảo dấu bằng chứng của bộ ba yếu hơn ngẫu nhiên, và ổn định khi lớp dương hiếm (6,6\%).
  \item Ước lượng chéo theo bệnh nhân: nếu ước lượng $\operatorname{AUC}_t$ trên chính các luật đã xây FKG, mỗi ảnh
        khớp với luật của chính nó và $\operatorname{AUC}_t$ bị thổi phồng.
  \item Chỉ dùng ảnh gốc: một mẫu SMOTE có thể có ảnh cha nằm ở phần khác, gây rò rỉ giữa các phần.
  \item Chi phí: thêm $K$ lần xây FKG trên tập huấn luyện (một lần cho mỗi nhóm dữ liệu và mỗi phép mờ hoá);
        suy diễn có cùng độ phức tạp với FISA-GM, $\mathcal{O}(|\mathcal{T}|)$ phép tra bảng và một tổng có trọng số.
  \item Trong M6 (FKGS), mô hình rút gọn dùng chung $\hat{\mathbf w}$ của tập huấn luyện đầy đủ; trong M3 và EPS,
        $\hat{\mathbf w}$ được ước lượng lại cho từng số mức và từng $\varepsilon$.
\end{itemize}
\end{nx}

\begin{algorithm}[htbp]
\caption{FISA-GM-w}\label{alg:mm.fisagmw}
\KwIn{Ảnh gốc huấn luyện đã mờ hoá $\{(\mathbf{x}_i,y_i,g_i)\}_{i=1}^{n_0}$; cơ sở luật huấn luyện đầy đủ
      (gồm mẫu SMOTE); truy vấn $\mathbf{x}$; $K=5$; $\varepsilon$.}
\KwOut{Nhãn $\hat y(\mathbf{x})$ và điểm $s_{\mathrm w}(\mathbf{x})$.}
Chia $\{1,\dots,n_0\}$ thành $P_1,\dots,P_K$ theo bệnh nhân, phân tầng theo nhãn\;
\For{$k=1,\dots,K$}{
  Xây FKG$^{(-k)}$ từ các ảnh gốc ngoài $P_k$\;
  \lForEach{$i\in P_k$, $t\in\mathcal{T}$}{$\tilde s_t(\mathbf{x}_i)\gets s^{(-k)}_t(\mathbf{x}_i)$ theo \eqref{eq:mm.triplescore}}
}
\lForEach{$t\in\mathcal{T}$}{$w_t\gets\max(\operatorname{AUC}_t-\tfrac12,0)$; chuẩn hoá $\hat w_t$ theo \eqref{eq:mm.tripleweight}}
Xây FKG từ toàn bộ cơ sở luật huấn luyện\;
$s_{\mathrm w}(\mathbf{x})\gets\sum_t \hat w_t\, s_t(\mathbf{x})$;\quad $\hat y(\mathbf{x})\gets\mathbb{1}[s_{\mathrm w}(\mathbf{x})>0]$\;
\end{algorithm}

% ---------------------------------------------------------------------------------------------------
% Câu đề xuất cho phần thực nghiệm (minh bạch về thời điểm đưa vào phương pháp):
% "FISA-GM-w được đưa vào sau khi kiểm toán kết quả ban đầu cho thấy FISA-GM gộp đều 560 bộ ba của Φsel,
%  trong đó khoảng 40% chứa từ hai đặc trưng ảnh có khả năng phân biệt thấp (AUC theo bộ ba 0,58–0,81), làm
%  loãng bằng chứng từ các thuộc tính lâm sàng. Trọng số được ước lượng chỉ trên ảnh gốc của tập huấn luyện,
%  áp dụng như nhau cho mọi mô hình FKG (kể cả FKG-UM), và không dùng tập kiểm tra để lựa chọn. Kết quả của
%  FISA-GM (không trọng số) được giữ lại ở Bảng ... để đối chiếu."
% ---------------------------------------------------------------------------------------------------
